\clearpage
\fancyhead[L]{\small ZACH. L'ATTRACTION DES MONTAGNES.}
\hypertarget{native-574-lower}{}
\begin{center}\small Göttingische gelehrte Anzeigen. Stück 210. Seite 2089..2100. 1814 December 31.\\[3pt]\rule{.72\linewidth}{.3pt}\end{center}

\emph{L'attraction des montagnes, et ses effets sur les fils à plomb ou sur les niveaux des
instrumens d'astronomie, constatés et determinés par des observations astronomiques et
géodésiques faites en} 1810, \emph{à l'ermitage de notre dame des anges, sur le mont de Mimet, et
au fanal de l'isle de Planier près de Marseille; suivis de la description géométrique de la
ville de Marseille et de son territoire. Par le baron} \textsc{de Zach}. 1814.
Avignon. Bei \textsc{Seguin} dem Aeltern. XX und 716 Seiten in gross Octav.

Die von \textsc{Bouguer} und \textsc{Condamine} im Jahre 1738 am Chimborasso, und von
\textsc{Maskelyne} im Jahre 1774 am Shehallien
 in Schottland über den Einfluss grosser Gebirgsmassen auf die Richtung der Schwere
angestellten Beobachtungen waren bis jetzt die einzigen directen Versuche in Beziehung auf
diesen interessanten Gegenstand gewesen. Die erstern hatten für die Ablenkung des Loths von
seiner natürlichen Richtung $7''5$, die andern $5''8$ gegeben. Allein eigentlich haben nur
die Beobachtungen in Schottland eine entschiedene Zuverlässigkeit, keinesweges aber die
unter zu ungünstigen Umständen und mit verhältnissmässig zu unvollkommenen Werkzeugen
ausgeführten Beobachtungen in Peru. Die Wichtigkeit des Gegenstandes macht es daher höchst
wünschenswerth, dass diese interessanten Erfahrungen bei dem gegenwärtig so sehr
vervollkommneten Zustande der practischen Astronomie mehr vervielfältigt werden mögen, und
die vorliegende Arbeit eines der ersten Beobachter ist demnach als ein wahrer Gewinn für die
Wissenschaft zu betrachten. Das Werk, welchem diese Anzeige gewidmet ist, enthält die
sämmtlichen zu dieser Operation gehörigen Beobachtungen bis ins kleinste Detail, entwickelt
ausführlich die angewandten Beobachtungs- und Rechnungsmethoden, und gibt ausserdem eine
Menge mit Sorgfalt berechneter Hülfstafeln, die auch sonst mit grossem Vortheil benutzt
werden können. Dass es überdies noch einen Reichthum an mancherlei practischen Bemerkungen
und Urtheilen, so wie an gelehrten Notizen und Nebenuntersuchungen enthält, ist ein
Vorzug, den man an den Schriften des Verfassers schon gewohnt ist.

\SourceLeaf{575}In der Einleitung führt der Verf. den Leser zuvörderst auf den Standpunkt, von welchem das Phänomen der Anziehung der Berge und die Beobachtungen, wodurch es sichtbar wird, betrachtet werden müssen, erzählt dann die darauf Bezug habenden bisher gemachten Erfahrungen, und gibt endlich eine allgemeine Uebersicht von seinen eignen Operationen. Hier mag es genug sein, in Erinnerung zu bringen, dass da die Schwere nur die Gesammtwirkung ist, welche alle Bestandtheilchen des Erdkörpers nach dem allgemeinen Anziehungsgesetze auf die an der Oberfläche der Erde befindlichen Körper ausüben, modificirt durch die aus der Rotation der Erde entstehende Centrifugalkraft, die astronomische Polhöhe, als der Winkel der Richtung der Schwere gegen die Ebene des Aequators, beim Aufsteigen vom Aequator nach dem Pole regelmässig und nach Gesetzen, die aus Vergleichungen der verschiedenen Gradmessungen abgeleitet werden, zunehmen muss, in so fern die Erde als ein regelmässiger Körper betrachtet wird; dass aber dies regelmässige Fortschreiten gestört wird da, wo grosse Abweichungen der Oberfläche der Erde von der Normalgestalt merkliche Ablenkungen der Schwere von der natürlichen Richtung hervorbringen. Der Unterschied der Polhöhen zweier Oerter, an deren einem die Richtung der Schwere von Norden nach Süden oder von Süden nach Norden durch eine nahe Bergmasse afficirt wird, wo hingegen an dem andern der Einfluss derselben entweder nicht mehr merklich ist, oder seine Wirkung in entgegengesetzter Richtung äussert, wird durch astronomische Beobachtungen anders gefunden werden, als durch die Rechnung aus der gegenseitigen Lage und Entfernung beider Oerter, und diese Verschiedenheit gibt uns die Grösse der Ablenkung im erstern Fall, oder die Summe beider Ablenkungen in andern zu erkennen. Offenbar ist also der zweite Fall, welcher bei \textsc{Maskelyne}'s Beobachtungen eintrat, der vortheilhaftere, und im erstern Fall das Geschäft doppelt schwieriger, indem die ganze Ablenkung immer nur sehr wenige Secunden beträgt, und die Beobachtungen daher von der grössten Feinheit sein müssen. Das Locale, wo Herr \textsc{von Zach} seine Operation ausführte, verstattete wahrscheinlich die Anwendung jenes vortheilhaftern Verfahrens nicht. Seine beiden Beobachtungspunkte waren, der eine am südlichen Abhange des Berges Mimet, des höchsten (400 Toisen über der Meeresfläche) in einer von Ost nach West laufenden Reihe von Kalkbergen, etwa zwei Meilen nördlich von Marseille, der andre der Leuchtthurm auf der sehr kleinen Insel Planier, ein Paar Meilen S. W. von Marseille. An dem erstern Orte, in der Höhe von 250 Toisen bei einem verfallenen Kloster Notre-Dame des Anges, musste der Berg eine Ablenkung der Richtung der Schwere nach Norden, folglich eine verminderte Polhöhe, bewirken: an dem andern Orte könnte der Einfluss der Berganziehung, wegen der grossen Entfernung, als unmerklich angesehen werden. Die beobachteten astronomischen Polhöhen mussten also einen kleinern Unterschied geben, als die unter Voraussetzung des regelmässigen Fortschreitens geführte Rechnung aus der durch geodätische Messungen bestimmten Lage beider Oerter. Zur Bestimmung der Polhöhen wendet man gern an beiden Oertern dieselben Sterne an; der Unterschied der Polhöhen, auf welchen allein es hier ankommt, wird dadurch unabhängig von der absoluten Richtigkeit der Declinationen der Sterne: man kann selbst die absoluten Polhöhen aus dem Spiele lassen, und anstatt ihres Unterschiedes sich an die Differenz der beobachteten Zenithdistanzen der Sterne halten, nachdem man sie von den kleinen periodischen Ungleichheiten befreiet und auf einerlei Zeitpunkt reducirt hat.

Von den acht Abschnitten, in welche Herr \textsc{von Zach} das vorliegende Werk getheilt hat, enthalten die beiden ersten die sämmtlichen bei Notre-Dame des Anges und auf der Insel Planier angestellten astronomischen Beobachtungen. Diese waren von dreierlei Art: Beobachtungen der Zenithdistanzen der drei Sterne $\alpha$ Ophiuchus, $\zeta$ Adler und $\alpha$ Adler; Beobachtungen von Pulversignalen zur Bestimmung der Längenunterschiede mit der Marseiller Sternwarte; und Beobachtungen von Azimuthen zur Orientirung des Dreiecksnetzes, wodurch die beiden Beobachtungsorte verbunden wurden. Die Zeitbestimmungen \SourceLeaf{576}geschahen an drei \textsc{Emery}'schen Chronometern mit Hülfe der correspondirenden Sonnenhöhen; die Zenithdistanzen wurden mit einem zwölfzolligen Vervielfältigungskreise von \textsc{Reichenbach} beobachtet, die Azimuthe mit einem achtzolligen Theodolithen von demselben Künstler. Schätzbar für die Beobachter mit Vervielfältigungskreisen ist die hier mitgetheilte allgemeine Reductionstafel für die ausser der Culmination beobachteten Zenithdistanzen, schärfer und vollständiger berechnet, als man sie anderswo findet. Der dritte Abschnitt enthält die terrestrischen Messungen. Eine Basis von 1182,4 Toisen wurde auf der Strasse von Marseille nach Aix mit hölzernen Messstangen gemessen, über deren geringe Veränderlichkeit der Verf. hier merkwürdige Erfahrungen beibringt. Die horizontalen Winkel der sieben Dreiecke des Netzes, das die beiden Hauptpunkte mit der Basis verband, wurden gleichfalls mit dem Theodolithen beobachtet, und meistens zehnmal, einige öfter, repetirt; die grösste Abweichung der Summe der drei Winkel in einem Dreieck war $4''9$. Im vierten Abschnitt leitet der Verf. aus diesen geodätischen Messungen, indem er die sphäroidische Gestalt der Erde, die Abplattung $\frac{1}{310}$, und den Halbmesser des Erdäquators 3271604 Toisen zum Grunde legt, nach den von \textsc{Delambre} in den Méthodes analytiques pour la détermination d'un arc du méridien gegebnen Formeln, den Unterschied der Breite und Länge für die beiden Beobachtungsorte ab: jener findet sich $12'\,3''11$; dieser $15'\,46''46$. Auch für diese Rechnungen theilt der Verf. verschiedene Hülfstafeln mit. Diese Rechnungsresultate werden nun im fünften Abschnitt mit den Resultaten der astronomischen Beobachtungen verglichen. Für den Unterschied der Polhöhen gaben

\begin{center}\small\begin{tabular}{@{}l@{\hspace{9pt}}r@{}}
598 Beobachtungen von $\alpha$ Ophiuchus $.\;.\;.\;.$ & $12'\,0''84$ \\
518 Beobachtungen von $\zeta$ Adler $.\;.\;.\;.\;.$ & $12\phantom{'}\,1.26$ \\
654 Beobachtungen von $\alpha$ Adler $.\;.\;.\;.\;.$ & $12\phantom{'}\,1.30$ \\
\end{tabular}\end{center}
\noindent also im Mittel 1770 Beobachtungen $12'\,1''13$, so dass $1''98$ als das letzte Resultat für die Wirkung der Anziehung des Mimet an dem Beobachtungsplatze Notre-Dame des Anges zu betrachten ist. Durch die Beobachtungen der Pulversignale hatte sich der Längenunterschied gefunden zwischen der Marseiller Sternwarte und

\begin{center}\small\begin{tabular}{@{}l@{\hspace{9pt}}r@{}}
Notre-Dame des Anges $.\;.\;.\;.\;.\;.\;.\;.$ & $7'\,29''25$ \\
Leuchtthurm auf der Insel Planier $.\;.\;.\;.$ & $8\phantom{'}\,5.70$ \\
\end{tabular}\end{center}
\noindent ersterer aus 63 Beobachtungen von 11 verschiedenen Tagen, letzterer aus 53 Beobachtungen von 12 Tagen, (wobei fünf Beobachtungen eines Tages ohne andern Grund, als weil sie ein von den übrigen zu abweichendes Resultat gaben, $7'\,38''1$, ausgeschlossen waren). Der ganze Längenunterschied zwischen Notre-Dame des Anges und der Insel Planier wird also aus diesen Beobachtungen $15'\,34''95$, oder auf die Punkte reducirt, auf welche sich die geodätischen Messungen beziehen, $15'\,35''79$, also um $10''67$, oder wenn man gar keine Beobachtung ausschliesst, um $13''05$ kleiner, als durch die geodätischen Messungen. Allein diesen Unterschied ist man keinesweges berechtigt, auch der Berganziehung zuzuschreiben, sondern vielmehr, wenigstens grösstentheils, den unvermeidlichen Beobachtungsfehlern der astronomischen Bestimmung des Längenunterschiedes, die, abhängig von der Zeit, bei weitem nicht des Grades von Genauigkeit fähig war, wie die Beobachtung der Zenithdistanzen. Uebrigens war die ganze Längenbestimmung nur eine untergeordnete, zum Hauptgeschäft gar nicht wesentliche Operation, und es ist schon genug, dass wir dadurch belehrt werden, welchen Fehlern auch der geübteste Beobachter bei dem angewandten Verfahren noch ausgesetzt bleibt. Eine ähnliche Belehrung geben die Vergleichungen der verschiedenen beobachteten Azimuthe. Wenn die auf der Insel Planier gemachten Azimuthalbestimmungen auf den Punkt Notre-Dame des Anges übertragen werden, so ergeben sich zwischen \SourceLeaf{577}den so abgeleiteten und den daselbst unmittelbar beobachteten Azimuthen Unterschiede von 7 und 21
Secunden, Abweichungen, welche durch eine billige Vertheilung auf die Azimuthalbestimmungen und die
terrestrischen Winkelbeobachtungen füglich erklärt werden können. Wir können jedoch nicht umhin
hier einen kleinen Umstand zu erwähnen, der gerade bei Vergleichungen \emph{dieser} Art nicht
unwichtig ist. Es ist uns aufgefallen, dass beim Centriren der Winkel die Distanzen der
Beobachtungsplätze von den Dreieckspunkten, welche der Verf. in Zehntausendtheilen von Toisen
angegeben hat, wenn sie auf metrisches Mass reducirt werden, fast sämmtlich nur Zehntheile vom
Meter, also runde Decimeter geben. Dies scheint doch kein Zufall zu sein, sondern der Verf. scheint
jene Distanzen mit einem metrischen Massstabe gemessen und kleinere Theile unbeachtet gelassen zu
haben. Inzwischen ändert bei mehrern Winkeln ein Fehler von Einem Decimeter (etwa 4 Zoll) die
Reduction um $3''$. Bei dem Hauptzweck des Verf. ist dies freilich etwas sehr unbedeutendes: allein
wenn man die Uebertragung eines Azimuths auf einen andern entfernten Orts und die Vergleichung mit
einem am letztern unmittelbar beobachteten benutzen will, um über die Gestalt der Erde neue
Aufschlüsse zu erhalten, so wird man auch bei diesem Geschäft des Centrirens eine desto grössere
Sorgfalt anwenden müssen, je grösser die Anzahl der Zwischen-Dreiecke ist. Der Verf. selbst
empfiehlt mit Wärme, die Ausführung einer solchen Operation mit \textsc{Reichenbach}'schen
Instrumenten, und schlägt dazu die Liparischen Inseln vor. Wir wünschen nichts mehr, als diese
Unternehmung von dem Verf. selbst ausgeführt zu sehen, wozu wir bei der Veränderung seines
bisherigen Aufenthalts um so mehr Hoffnung haben dürfen.

Der sechste Abschnitt ist der gehaltreichste des ganzen Werks. Der Verf. untersucht
zuvörderst den Einfluss, welchen die möglichen Fehler der einzelnen Operationen auf das
Endresultat haben können. Man begreift leicht, dass die Beobachtungen der Zenithdistanzen bei
weitem der delicateste Theil des Ganzen sind, und in der That könnte mancher bei der
ausserordentlichen Kleinheit des Endresultats, dessen Zuverlässigkeit eine Genauigkeit selbst bis
auf Theile von Secunden bei jenen voraussetzt, um so eher Bedenklichkeit haben, da der Verfasser
selbst früher bei anderer Gelegenheit erklärt hat, dass man bei allen mit Hülfe von
Repetitionskreisen bestimmten absoluten Polhöhen immer einer Ungewissheit von mehrern Secunden
ausgesetzt bleibe. Freilich gründete sich diese Behauptung auf die Erfahrung, dass solche
Bestimmungen durch \emph{verschiedene} Kreise gemacht dergleichen Unterschiede zeigten, während die
Resultate durch jeden einzelnen Kreis unter sich vortrefflich übereinstimmten. Bei der
gegenwärtigen Operation hingegen sind alle Beobachtungen mit einem und demselben Instrumente
gemacht, sie haben unter sich die schönste Uebereinstimmung, und es kommt hier nicht auf die
absoluten Bestimmungen selbst, sondern nur auf kleine Unterschiede an. Herr \textsc{von Zach}
benutzt zu einer weitern Bestätigung seiner Resultate die zahlreichen astronomischen
Beobachtungen, welche er theils auf der Königl. Sternwarte zu Marseille, theils auf seinen eignen
Sternwarten La Capellette und St. Peyre angestellt hat. Die Unterschiede zwischen den dadurch
bestimmten Polhöhen dieser drei Punkte (welche von der Gebirgsanziehung nicht mehr merklich
afficirt werden) 
und der Polhöhe der Insel Planier stimmen äusserst nahe mit den Resultaten der geodätischen
Messungen überein, wodurch er die Verbindung mit seinem Dreiecknetze bewirkt hat: während eine
ähnliche Vergleichung mit Notre-Dame des Anges wieder ziemlich übereinstimmend jenen Unterschied
von zwei Secunden herbeiführt. Man könnte die Zulässigkeit dieses Bestätigungsgrundes vielleicht in
Zweifel ziehen, in so fern sich diese fünf Polhöhen auf verschiedene Sterne gründen (Notre-Dame des
Anges und die Insel Planier auf $\alpha$ Adler, hingegen die drei andern auf $\alpha$ und $\beta$
im kleinen Bär). Allein wenn man gehörig erwägt, dass die dabei angewandten Declinationen gerade
diejenigen sind, welche Herr \textsc{von Zach} selbst mit eben diesem Instrumente an einem nicht
sehr viel nordlicher liegenden Orte, Mailand, bestimmt hat, so überzeugt man sich leicht von der
Nichtigkeit dieses Zweifels, in so fern man nur die Voraussetzung gelten lässt, dass die \SourceLeaf{578}etwaigen Fehler, die bei den mit dem Kreise bestimmten absoluten Zenithdistanzen Statt finden mögen, für bestimmte Zenithdistanzen unveränderlich, und für wenig verschiedene Zenithdistanzen sehr nahe gleich sind. Was man auch immer von den Ursachen der merkwürdigen von Hrn. \textsc{von Zach} mit \emph{verschiedenen} Kreisen gefundenen Differenzen urtheilen mag, so ist doch nicht zu leugnen, dass ohne die erwähnte Voraussetzung die bleibende Harmonie der mit einerlei Kreise gefundenen Resultate sich nicht erklären lasse. Nach allen diesen Gründen darf man annehmen, dass in der That der gefundene Unterschied von 2 Secunden sehr nahe die Anziehung des Berges Mimet darstellt, dabei wird jedoch kein praktischer Astronom in Abrede sein, dass in diesem Resultate immer noch eine Ungewissheit von einer halben Secunde (wo nicht mehr) zurückbleibe, und sich also dasselbe in keine engeren Grenzen als $1''5$ bis $2''5$ einschränken lasse. Man muss daher allerdings bedauern, dass eine so schöne und so sorgfältig ausgeführte Operation nur eine so kleine Grösse hervorgebracht hat, gegen welche die unvermeidliche Ungewissheit in einem so bedeutenden Verhältnisse steht. Es ist also um so mehr Schade, dass der zweite Beobachtungsort nicht, anstatt auf der Insel Planier, auf der Nordseite unmittelbar am Berge Mimet genommen ist. Herr \textsc{von Zach} erklärt sich nicht über die Gründe, welche ihn davon abgehalten haben: allein ein so einsichtsvoller Astronom hätte gewiss nicht auf einen Vortheil, der die Wirkung vielleicht verdoppelt hätte, Verzicht geleistet, wenn nicht das Local unübersteigliche Hindernisse dargeboten hätte. Wäre die Wirkung des Mimet bedeutender ausgefallen, so würden wir auch noch sehr eine vollständigere Kenntniss von den körperlichen Abmessungen dieses Berges gewünscht haben: allein unter den obwaltenden Umständen würden doch die Schlüsse, welche man daraus auf die comparative Dichtigkeit des Minet und des ganzen Erdkörpers machen könnte, eine zu beschränkte Genauigkeit geben. Immer aber gewähren die \textsc{von Zach}'schen Messungen, indem sie wenigstens die ausserordentliche Kleinheit des Einflusses einer so bedeutenden Bergmasse beweisen, den wichtigen Nutzen, dass sie uns gegen eine zu voreilige Berufung auf den möglichen Einfluss von Localattractionen, wenn die Messungen nicht zusammenpassen wollen, etwas misstrauischer machen.

Ausser den angezeigten Untersuchungen enthält der sechste Abschnitt ferner die Prüfung einiger Bestimmungen von \textsc{Cassini de Thury} in der Méridienne vérifiée, die von denen des Verf. merklich abweichen. Noch viel grössere und in der That bis zur Entstellung gehende Fehler finden sich in \textsc{Cassini}'s description géometrique de la France und in seiner grossen Charte von Frankreich. --- Endlich gibt der Verf. noch eine Uebersicht der Resultate, welche für die Stellungen verschiedener Sterne theils aus seiner gegenwärtigen Arbeit, theils aus seinen eignen frühern Beobachtungen, theils aus einer kritischen Discussion der Beobachtungen anderer Astronomen folgen. Wir erwähnen davon hier nur die in vielfacher Beziehung so wichtige gerade Aufsteigung des Polarsterns, für welche \textsc{Mathieu}'s Beobachtungen im Jahre 1812 eine Vermehrung von 5 Zeitsecunden gegeben haben. Um hiemit \textsc{von Zach}'s Bestimmung von 1790 in Uebereinstimmung zu bringen, muss man eine eigne Bewegung voraussetzen, und diese findet Herr \textsc{von Zach} durch die Vergleichung mit \textsc{Lacaille}'s Beobachtungen von 1750 auch vollkommen bestätigt, und setzt sie auf $+3''177$ in Bogen jährlich. (Die schärfste Bestimmung der gegenwärtigen Rectascension des Polarsterns werden uns die zahlreichen seit fünf Jahren auf der Seeberger Sternwarte angestellten Beobachtungen des Hrn. \textsc{von Lindenau} geben, wovon wir schon vorläufig sagen können, dass sie die von \textsc{Mathieu} gefundene Vergrösserung bestätigen.)

Der beschränkte Raum unserer Blätter erlaubt uns den Inhalt der übrigen Abschnitte des \textsc{von Zach}'schen Werks nur noch kurz zu berühren. Der siebente Abschnitt bestimmt die Höhe der Dreieckspunkte und einiger anderer Oerter über dem mittelländischen Meere. Der achte Abschnitt enthält die Lage einer grossen Anzahl von Punkten in der Stadt Marseille und der umliegenden Gegend nach ihrem Abstande vom Meridian und Perpendikel der dortigen Königl. Sternwarte und nach ihrer Länge \SourceLeaf{579}und Breite, und überdies noch mehrere andere interessante kritische Untersuchungen über
verschiedene Punkte in Marseille, welche in der Geschichte der Astronomie merkwürdig geworden
sind. Endlich gibt Herr \textsc{von Zach} in einem Anhange noch eine neue Berechnung der
sämmtlichen von \textsc{Maskelyne} am Shehallien angestellten astronomischen Beobachtungen,
wodurch indessen das von \textsc{Maskelyne} selbst gefundene Endresultat keine Veränderung
leidet.
